% /root/monodromy, 2026-11-19. Proved; independent root and reduction_descent audits passed.
\subsection{Uniform Hasse geometry for every ordinary elliptic product}
\label{subsec:ordinary-products-general-structure}
\status{Proved and checked. The parent and reduction/descent agent independently audited the complete proof. The universal polynomial symmetrizer calculation was independently rerun by both agents. Full group exclusion is treated separately.}

\begin{lemma}[A universal elliptic Cartier symmetrizer]
\label{lem:ordinary-elliptic-universal-symmetrizer}
Let $E:y^2+xy=x^3+A x^2+B$ over a perfect field of characteristic
two, with $B\ne0$, and put $N=\mathcal O_E(3O)$.
In the basis $(1,x,y)$ and with the Cartier-invariant differential
$\eta=dx/x$, the root-removed Cartier multiplication matrix is
\[
 M(s)=\begin{pmatrix}
 s_0&B s_2&B s_1\\
 s_2&s_1+A s_2&s_0+A s_1\\
 0&s_2&s_1+s_2
 \end{pmatrix}.
\]
The constant invertible matrix
\[
 R=\begin{pmatrix}1&0&B\\0&0&B\\0&B&(1+A)B\end{pmatrix},
 \qquad\det R=B^2,
\]
satisfies
\[
 RM(s)=\begin{pmatrix}
 s_0&0&B s_2\\
 0&B s_2&B(s_1+s_2)\\
 B s_2&B(s_1+s_2)&B s_0+B s_1+(1+A)B s_2
 \end{pmatrix}.
\]
Thus $RM(s)$ is symmetric and its diagonal is $R^t s$.
For any pair of ordinary elliptic curves with degree-three
line bundles, the analogous nine-dimensional Cartier pencil
has a constant invertible symmetrizer and an invertible
linear map to its nine diagonal entries.
\end{lemma}

\begin{proof}
The smooth ordinary Weierstrass equation can be put in the
displayed form by first normalizing its nonzero $xy$
coefficient, then translating $x$ to remove the $y$ term and
translating $y$ by a constant to remove the $x$ term.
These changes preserve $O$ and the degree-three line bundle.
The discriminant in the resulting equation is $B$, so $B\ne0$.
A degree-three line bundle is a translate of $\mathcal O(3O)$
over the algebraic closure, which suffices for the geometric
assertion.

Writing a rational function as $a^2+x b^2$ gives
$\mathcal C((a^2+x b^2)\eta)=a\eta$.
The Weierstrass relation consequently gives
\[
 \mathcal C(1)=1,\quad \mathcal C(x)=0,\quad
 \mathcal C(y)=x,\quad
 \mathcal C(xy)=\sqrt B+\sqrt A\,x+y.
\]
Together with $\mathcal C(x^2)=x$ and
$\mathcal C(y^2)=y$, this is precisely the matrix written above
when the output scalar square root is removed.
Direct multiplication proves the formulas for $R$.
The matrix $R$ itself need not be symmetric; only $RM(s)$ is.

On a product the Cartier multiplication coefficients are tensor
products of the two one-factor coefficient matrices. Therefore
$R_1\otimes R_2$ symmetrizes all nine coefficient matrices.
Its diagonal map is $R_1^t\otimes R_2^t$, and its determinant
is $B_1^6B_2^6\ne0$.
The universal polynomial identities over
$\mathbf F_2[A_1,B_1,A_2,B_2]$ are independently reproducible in
\path{agents/monodromy/code/ordinary_product_universal_symmetrizer.py};
the complete retained output is
\path{agents/monodromy/outputs/ordinary_product_universal_symmetrizer.json}.
\end{proof}

\begin{lemma}[Squarefreeness from independent symmetric diagonals]
\label{lem:symmetric-independent-diagonal-squarefree}
Let $S(z)$ be an $m\times m$ symmetric matrix of linear forms
in characteristic two, and suppose its diagonal entries are
independent linear coordinates $z_1,\ldots,z_m$.
Then $\det S$ is geometrically squarefree.
\end{lemma}

\begin{proof}
In the determinant expansion, permutations having a cycle of
length at least three cancel with the inverse of the first such
cycle. The surviving terms are indexed by matchings:
\[
 \det S=\sum_{\mathcal M}
 \prod_{\{i,j\}\in\mathcal M}S_{ij}^2
 \prod_{i\text{ unmatched}}z_i.
\]
Only the empty matching contributes to the monomial
$z_1\cdots z_m$, and its coefficient is one. Every term of a
nonempty matching includes a square of a linear form and hence
has some exponent at least two in each of its monomials.
If a nonconstant polynomial $q$ had $q^2\mid\det S$ over an
algebraic closure, write $\det S=q^2r$ with $\deg r<m$.
Every exponent in $q^2$ is even. A term contributing to
$z_1\cdots z_m$ would therefore require a monomial of $r$
with all $m$ exponents odd, whose degree is at least $m$.
This is impossible. Thus no irreducible factor is repeated.
\end{proof}

\begin{lemma}[Jet bounds for arbitrary ordinary cubic factors]
\label{lem:ordinary-product-uniform-jets}
Let $X=E_1\times E_2$ with both elliptic curves ordinary in
characteristic two and $L=\mathcal O_{E_1}(3O)\boxtimes
\mathcal O_{E_2}(3O)$. The conclusions of
Lemma~\ref{lem:product-small-degree-jets} hold on $X$.
In addition, the singular-curve discriminant of $|L|$ is
irreducible with a general one-node member, and its general
member is ordinary.
\end{lemma}

\begin{proof}
The tensor first-jet proof of that lemma uses only the plane
cubic embeddings. Its one finiteness point is also uniform:
for $F=Y^2Z+XYZ+X^3+A X^2Z+B Z^3$ the partials are
\[
 F_X=YZ+X^2,\qquad F_Y=XZ,\qquad
 F_Z=Y^2+XY+A X^2+B Z^2.
\]
Their constant linear combinations are linearly independent,
by comparing $Y^2,X^2,XZ$. No nonzero quadratic vanishes
identically on the cubic, so the tangents through any fixed
point form a finite set. Thus the exceptional rank-five pairs
are the pairs sharing one elliptic coordinate, of dimension
three in $X\times X$, together with the crossed incidence
locus, which is finite over the first point and has dimension
at most two. Elsewhere the two first jets have rank six.

A specified singular point imposes three independent conditions,
so its incidence is an irreducible projective bundle over $X$.
An ordinary node occurs in a nonempty open of every fibre:
products of sections vanishing simply in the two elliptic
directions provide an arbitrary nonzero mixed quadratic
coefficient. For two singular points, the rank-six locus gives
codimension $6-4=2$ after the points vary. The exceptional
rank-five loci give codimensions at least $5-3=2$ and $5-2=3$.
Hence the general discriminant member has exactly one node,
and nonreduced members have codimension at least two.

Finally $|L|$ contains an ordinary grid: take three distinct
vertical fibres with sum $3O$ on $E_1$ and three distinct
horizontal fibres with sum $3O$ on $E_2$.
Its normalization components are ordinary elliptic curves and
its graph term is Frobenius-bijective, so the grid is ordinary.
The irreducible discriminant thus contains an ordinary point,
and its ordinary open is nonempty and dense.
\end{proof}

\begin{theorem}[Uniform irreducible Hasse divisors on ordinary products]
\label{thm:ordinary-products-Hasse-irreducible}
Let $E_1,E_2$ be arbitrary ordinary elliptic curves over a
perfect field of characteristic two, and let
$L=\mathcal O_{E_1}(3O)\boxtimes\mathcal O_{E_2}(3O)$.
The full Hasse divisor in $|L|$ is geometrically reduced and
has generic Cartier corank one on every component.
For every $d\ge2$, the full Hasse divisor in $|L^d|$ is
geometrically irreducible and reduced, with generic stable
rank loss exactly one. Consequently the geometric relative
Hasse--Witt group in degree $d\ge2$ is generated by
transvections.
\end{theorem}

\begin{proof}
The two preceding algebraic lemmas make the degree-one Hasse
polynomial geometrically squarefree. Its derivative is nonzero
at each generic component point, so its matrix has corank one
there. By Lemma~\ref{lem:ordinary-product-uniform-jets}, no
component is the singular-curve discriminant, whose generic
point is ordinary. Thus the generic curve on each degree-one
Hasse component is smooth with a one-dimensional Cartier
kernel. Stable loss one in this degree is not assumed or
needed in the following argument.

Fix $d\ge2$, put $e=d-1$, and repeat the mixed-tacnode proof
with $C\in|L|$ and a general ordinary $D\in|L^e|$.
We spell out the differences caused by using $|L|$ in place
of $|L^3|$.
For a smooth general $C$ on a base Hasse component, the
nonzero tangent functional $dh_1$ belongs to the dual
first-jet space at only finitely many points of $C$.
Indeed, if there were infinitely many, choose one such $P$;
all but finitely many other points $Q$ on $C$ avoid the two
fibres and the finite crossed set through $P$.
The first-jet dual spaces at $P$ and such a $Q$ are disjoint,
contradicting that the same nonzero $dh_1$ lies in both.
We can therefore choose $P$ avoiding those finitely many
points, the zeros of the Cartier-killed differential, and the
flex fibres. The affine tangent cone $\ker dh_1$ maps
surjectively to first jets at this $P$.

The incidence with $D$ tangent to $C$ at $P$ consequently
projects dominantly to $|L^e|$, by the same three-equation
infinitesimal calculation in
Theorem~\ref{thm:product-Hasse-irred}. The good open of smooth
$D$ with exactly one simple contact is nonempty. For $e\ge2$
use full two-point first jets; for $e=1$ use rank at least five
and its exceptional one-dimensional locus through $P$.
The latter meets $C$ in only finitely many points. Exactly the
same dimension counts as in that theorem exclude a second
contact or a singularity of $D$.
Hence a single general ordinary $D$ admits such a mixed
contact with every base Hasse component.

For completeness the fixed-$D$ contact divisor is still
irreducible with generic ordinary $C$ despite the weaker
base jets. Its incidence over $P\in D$ is a projective bundle
cut out by two independent conditions. Two contacts impose
four conditions on the rank-six locus of pairs in $D^2$;
the rank-five exceptional pairs have dimension at most one
and impose at least three. Both give codimension at least
two in $|L|$. A singularity of $C$ at a second point imposes
five conditions generically, with the point pair varying in
dimension three; on the exceptional locus the rank is at
least four and the dimension is at most two. A singularity
at the contact point also has codimension at least two.
Order-three contact imposes three conditions at a nonflex
point; the flex fibres meet $D$ in finitely many points, so
even there the two contact conditions give codimension two.
The generic contact is therefore simple with $C$ smooth.
The joint tangency incidence for $(P,C,D)$ is irreducible
and dominates both complete linear systems by first jets.
Its two ordinary inverse images are nonempty dense opens,
so the generic contact $C$ is ordinary for a general ordinary
$D$.

On the multiplication subspace $g_D|L|$, the divisorial
nonordinary support is the union of the base Hasse components
and this contact divisor, by the exact union formula.
The exceptional nonreduced base locus has codimension at least
two; when $e=1$, $C=D$ is a single point of $\mathbf P^8$,
and when $e>1$ no base member contains $D$.
At each mixed contact, Lemmas~\ref{lem:pointed-Cartier} and
\ref{lem:mixed-tacnode-Hasse} make the full Hasse divisor
smooth. Their proofs use only ordinarity of the ambient product,
not the particular elliptic equation, and require only Cartier
corank one on $C$, not stable rank loss one on $C$. Therefore every base
Hasse component and the contact divisor are assigned to the
same global Hasse component. Every global component has a
positive-degree nonzero restriction to this linear subspace,
which contains an ordinary point. The component-merging
argument of Theorem~\ref{thm:product-Hasse-irred} proves that
there is exactly one global component. The smooth mixed
point also makes its multiplicity one.

There is an ordinary pair $(C,D)$ with a single simple
contact. The union formula gives stable loss one at that
point of the now irreducible Hasse divisor. The closed
iterated-Frobenius rank loci imply generic stable loss at
most one, and noninvertibility makes it exactly one.
Lemma~\ref{lem:char2-defect-one-inertia} proves the last assertion.
\end{proof}

\begin{corollary}[An irreducible kernel for every ordinary pair]
\label{cor:ordinary-products-irreducible-kernel}
In the preceding theorem let $d\ge2$, put $b=9d^2$, and set
$q=0$ for odd $d$ and $q=3$ for even $d$.
The entire invariant-covector space has dimension $q$.
Its common kernel $K\subset V_d$ is irreducible of dimension
$b-1-q$, and $V_d$ has no invariant vector.
The image on $K$ is irreducible and generated by transvections.
\end{corollary}

\begin{proof}
The fixed-covector calculation is already
Theorem~\ref{thm:product-char2-all-parities-fixed-quotient}.
Its statement and proof allow distinct ordinary factors: the
one-factor Frobenius pushforward calculations and their external
tensor product do not require the elliptic curves to be isomorphic.
The relative coherent bundle on $\mathbf P^{b-1}$ is
$\Omega^1_{\mathbf P^{b-1}}$, by the Euler/adjunction sequence
and the ordinary one-dimensional ambient $H^2(\mathcal O_X)$.
Theorem~\ref{thm:ordinary-products-Hasse-irreducible} therefore
permits Lemma~\ref{lem:irreducible-hasse-quotients-trivial}.
Every nonzero proper invariant subspace has trivial quotient,
so contains the common kernel. This makes $K$ irreducible and
rules out invariant lines, since $\dim K\ge32$.
The transvection assertion follows by restriction; a nontrivial
generator restricts either to a transvection or to the identity.
\end{proof}
